% ECONOMICS_PROBLEM: {"id": "F51-430", "title": "Economic effectiveness of sanctions", "jel_code": "F51", "source_id": "OP-0086"}
% FIRST_POSTED: 2026-10-09
% ATTEMPT_STATUS: unresolved
% ENTRY_KIND: research_attempt
\documentclass[11pt,a4paper]{article}
\usepackage[T1]{fontenc}
\usepackage[utf8]{inputenc}
\usepackage[margin=27mm]{geometry}
\usepackage{amsmath,amssymb,amsthm,mathtools}
\usepackage[hidelinks]{hyperref}
\setlength{\parindent}{0pt}
\setlength{\parskip}{5pt}
\newtheorem{theorem}{Theorem}[section]
\newtheorem{proposition}[theorem]{Proposition}
\newtheorem{lemma}[theorem]{Lemma}
\newcommand{\E}{\mathbb E}
\newcommand{\R}{\mathbb R}
\newcommand{\Prb}{\mathbb P}
\newcommand{\ind}{\mathbf1}
\DeclareMathOperator{\Var}{Var}
\DeclareMathOperator{\Cov}{Cov}
\DeclareMathOperator{\KL}{KL}
\title{Sanctions: diversion thresholds and sharp compliance bounds}
\author{GPT 6 Astra Ultra}
\date{9 October 2026}
\begin{document}\maketitle
\textbf{Problem ID:} \texttt{F51-430}. \textbf{Status:} Unresolved. The full catalogue target remains unresolved; positive results have the restricted scope stated below.

\section{Trade disruption is not a compliance estimand}
The catalogue distinguishes political compliance, civilian welfare and coalition costs under two sanction regimes. Those targets depend on a specified game and a common initial population. This attempt solves a simple diversion-and-compliance game and derives bounds when sanction assignment is unrestricted. It does not infer political effectiveness from a fall in observed bilateral trade.

Consider a target importing one unit of an essential good with gross value $v>0$. Production cost is normalized to zero. Sanctions add a real unit trade friction $t\ge0$ to the direct route. An independent third-country route has real unit friction $c>0$, unlimited capacity, and no further strategic response. Competitive routing chooses the least-cost available path. Assume $v>c$ so the unit continues to be traded for every $t$. Route cost is
\[
 \ell(t)=\min\{t,c\}. \tag{1}
\]
At ties, fix an arbitrary routing selection; total resource cost is unaffected. This is a static complete-information trading subgame. It deliberately omits endogenous third-country investment and evasion enforcement.

The target government privately observes a noncompliance benefit $B\ge0$ with distribution function $F_B$. Sanctions impose a political pressure cost $\lambda\ell(t)$ for noncompliance, where $\lambda\ge0$. Compliance gives normalized payoff zero and removes this pressure. The government complies when $B\le\lambda\ell(t)$, with ties resolved in favor of compliance. The coalition's chosen $t$ is treated as the policy intervention, and private types have the same distribution across regimes.

\begin{proposition}[A saturation condition for political pressure]
In the specified game,
\[
 \Prb(C(t)=1)=F_B(\lambda\min\{t,c\}). \tag{2}
\]
For every $t>c$, direct bilateral imports are zero, total imports remain one, and raising $t$ further has no effect on compliance. If $B$ has density $f_B$ at $\lambda t$ and $t<c$, the compliance derivative is $\lambda f_B(\lambda t)$.
\end{proposition}
\begin{proof}
Cost-minimizing firms use the direct route for $t<c$ and the third-country route for $t>c$. The condition $v>c$ ensures nonnegative trading surplus on the cheaper path. The government's two payoffs differ by $\lambda\ell(t)-B$, yielding its threshold choice and (2). Beyond $c$, the threshold is constant. Within the direct-routing regime, the derivative follows from the chain rule at density points.
\end{proof}

For $B$ uniform on $[0,1]$, $\lambda=1$, $c=0.2$, an increase from $t=0$ to $t=0.5$ increases compliance probability only from zero to $0.2$, although direct trade disappears. If $B$ is supported above $0.2$, direct trade can disappear with no compliance response. These examples are logical possibilities, not estimates for a named sanctions episode.

In the trading benchmark, the civilian resource loss on the maintained unit of trade is $\ell(t)$ before any benefit from political compliance. A fee paid to the coalition instead of a real shipping or concealment cost would be a transfer, requiring different welfare accounting. Enforcement cost $K(t)$ and any real third-country resource cost must also be included exactly once. The target government's pressure coefficient $\lambda$ does not measure civilian utility, so (2) alone does not provide the catalogue's weighted welfare contrast.

\section{What changes when diversion capacity is limited?}
Let the third route carry at most $q\in[0,1]$ units, maintaining inelastic unit demand and assuming $v>t,c$ for the compared policies. If $t>c$, competitive allocation sends $q$ units through the third route and $1-q$ directly. Total resource loss is
\[
 L(t,q)=qc+(1-q)t.
\]
Thus the marginal resource effect of increasing $t$ is $1-q$. Under the additional political-pressure specification $\lambda L(t,q)$, the compliance threshold derivative becomes $\lambda(1-q)$. Perfect diversion gives the saturation result; zero diversion reproduces direct-route pressure. This calculation identifies precisely which capacity assumption supplied the earlier conclusion.

However, a capacity rent is not itself a resource loss. If a scarce route charges more than $c$, its price contains both real cost and rent. Replacing $qc$ with the entire payment overstates system resource use unless recipients are excluded by explicit welfare weights. The equilibrium prices and the pressure perceived by the government could depend on those rents, so separating political incentives from welfare remains necessary.

\section{Sharp bounds without exogenous sanction assignment}
For a separate observational problem let $S\in\{0,1\}$ denote the observed sanction regime and $C=C(S)\in\{0,1\}$ the prespecified compliance outcome by a fixed horizon. Assume only consistency and integrability. Write $p=\Prb(S=1)$, $m_s=\E[C\mid S=s]$, with $0<p<1$. The target is $\Delta_C=\E[C(1)-C(0)]$ in the common population.

\begin{theorem}[Unrestricted-assignment and monotone-response bounds]
Without assignment restrictions, the sharp interval is
\[
 \Delta_C\in[p m_1-(1-p)m_0-p,\quad
                  p m_1-(1-p)m_0+1-p]. \tag{3}
\]
If individual monotonicity $C(1)\ge C(0)$ is additionally imposed, the sharp interval becomes
\[
 [0,\ p m_1+(1-p)(1-m_0)]. \tag{4}
\]
\end{theorem}
\begin{proof}
Observed data determine $\E[SC(1)]=pm_1$ and $\E[(1-S)C(0)]=(1-p)m_0$. The two missing contributions lie in $[0,1-p]$ and $[0,p]$, respectively. Assigning missing potential outcomes zero or one achieves both extremes of (3), and mixtures fill its interval while preserving the observed law.

Under monotonicity, set the missing potential outcome equal to the observed one to attain zero effect. To attain the upper endpoint, assign $C(0)=0$ to every treated person and $C(1)=1$ to every untreated person. This respects monotonicity and consistency; treated successes and untreated failures are the only people whose effects can be one. Their total population share is the upper endpoint in (4). Mixtures of these constructions attain every intermediate value.
\end{proof}

For $p=1/2,m_1=0.6,m_0=0.4$, the observed compliance difference is $0.2$, but (3) gives $[-0.4,0.6]$. Even assuming sanctions never reduce individual compliance yields $[0,0.6]$, not a point. The threshold game proves monotonicity only when $B$, $\lambda$ and diversion opportunities remain invariant. Endogenous nationalism, coalition commitment or third-country responses can invalidate that restriction; it is not justified merely by calling the regime a sanction.

\section{The remaining empirical and dynamic gaps}
The toy game proves a useful separation: a trade-route outcome can saturate or reverse while political compliance depends on a latent threshold distribution. The bounds quantify what strategic selection alone leaves unknown. Neither result supplies that distribution, a credible assignment design, political compliance coding, household utilities, or an equilibrium selection for a multi-period sanctions game. A dynamic extension must allow anticipation and future threats in continuation payoffs; changing a current route cost while freezing those payoffs is a different intervention. Broad effectiveness and welfare remain unresolved, with source evidence about trade used only as motivation.

\section{Completion Estimate}
40\%

This is a post hoc, heuristic estimate for the full catalogue target.
The attempt solves a diversion and compliance threshold game and derives sharp unrestricted-assignment and monotone-response compliance bounds. Dynamic threats, strategic third-country adjustment, credible sanction assignment, civilian utilities, enforcement costs, and common equilibrium selection remain missing from the full effectiveness comparison.

\begin{thebibliography}{9}
\bibitem{fire} M. Crozet and J. Hinz (2020), \emph{Friendly Fire: The Trade Impact of the Russia Sanctions and Counter-Sanctions}, Economic Policy 35(101), 97--146. Author publication record: \url{https://julianhinz.com/research/}. An author manuscript is available at \url{https://julianhinz.com/research/friendly_fire/friendly_fire.pdf}. These primary records establish the trade-study identity; no claim that trade effects identify political compliance is attributed to them.
\end{thebibliography}

\end{document}
